
\begin{description}

% \item[Geogebra] 

% \link{  .ggb}{GeoGebra applet}

\item[R] 

\link{http://www.math.unl.edu/~sdunbar1/markovchain.R}{R script for markovchain.}

\begin{lstlisting}[language=R]
library(markovchain)

stateNames <- c("R", "N", "S")
transMatrix <- matrix( c(1/2, 1/4, 1/4, 1/2, 0, 1/2, 1/4, 1/4, 1/2),
                      nrow=3, byrow=TRUE)
startState <- "N"

weatherOz <- new("markovchain", transitionMatrix=transMatrix,
                 states=stateNames, name="Weather_in_Oz")
print( weatherOz )

print( summary(weatherOz) )

pathLength <- 10 
weatherPath <- rmarkovchain(n=pathLength,
                            object = weatherOz, t0 = startState)
cat("A sample starting from N: ", weatherPath, "\n")
weatherPath <- rmarkovchain(n=pathLength,
                            object = weatherOz, t0 = startState)
cat("Another sample starting from N: ", weatherPath, "\n")

largeN <- 1000; startStable <- 100
weatherPath <- rmarkovchain(n=largeN,
                            object = weatherOz, t0 = startState)
stablePattern <-  weatherPath[startStable:largeN]
empiricalStable <- c( mean( stablePattern == "R" ),
                     mean( stablePattern == "N" ),
                     mean( stablePattern == "S" ))
theorStable <- steadyStates(weatherOz)
cat("Empirical stable distribution: ",  empiricalStable, "\n")
cat("Theoretical stable distribution: ", theorStable, "\n")
\end{lstlisting}  

% \item[Octave]

% \link{http://www.math.unl.edu/~sdunbar1/    .m}{Octave script for .}

% \begin{lstlisting}[language=Octave]

% \end{lstlisting}

% \item[Perl] 

% \link{http://www.math.unl.edu/~sdunbar1/    .pl}{Perl PDL script for .}

% \begin{lstlisting}[language=Perl]

% \end{lstlisting}

% \item[SciPy] 

% \link{http://www.math.unl.edu/~sdunbar1/    .py}{Scientific Python script for .}

% \begin{lstlisting}[language=Python]

% \end{lstlisting}

\end{description}



%%% Local Variables:
%%% mode: latex
%%% TeX-master: t
%%% End:
